\documentclass[11pt]{article}
\usepackage[a4paper,margin=25mm]{geometry}
\usepackage{amsmath,amssymb,amsthm,hyperref}
\newtheorem{theorem}{Theorem}\newtheorem{corollary}{Corollary}
\title{Orthogonal error spaces for two ordered product quadratures}
\author{Auxiliary exact identity; independently reviewed}\date{30 September 2026}
\begin{document}\maketitle
The sharp simultaneous two-die question remains open in this work.
This note reformulates its real hybrid version as orthogonality of
cumulative quadrature errors. It is not a new upper or lower size bound.

Let $\mu_1,\mu_2$ be probability measures on $[0,1]$ with no common
atoms, and put $\delta_i=\mu_i-dt$. For a polynomial $f$ define
\[
 E_{i,f}(t)=\int_{[0,t]}f(x)\,d\delta_i(x).
\]
Endpoint conventions do not affect the Lebesgue integrals below.

\begin{theorem}\label{main}
Suppose each $\mu_i$ integrates polynomials through degree $M+1$
exactly as Lebesgue measure, and the product rule integrates every
polynomial of total degree at most $M$ exactly on each open chamber
$x<y$ and $y<x$. Then
\begin{equation}\label{orthogonal}
 \int_0^1 E_{1,f}(t)E_{2,g}(t)p(t)dt=0
 \quad\text{if }\deg f+\deg g+\deg p\le M-1.
\end{equation}
In particular the two spaces
\[
 \mathcal E_i^{(r)}=\{E_{i,f}:\deg f\le r\}\subset L^2[0,1]
\]
are orthogonal whenever $2r\le M-1$. If $\mu_i$ is finitely supported,
each of these spaces has dimension exactly $r+1$.
\end{theorem}
\begin{proof}
For $\deg f+\deg g\le M$, expand the ordered product error as
\[
 C(f,g)=\int_{x<y}f(x)g(y)\,d\delta_1(x)d\delta_2(y).
\]
Its $\mu_1\otimes\mu_2$ and Lebesgue product terms are equal by the
chamber assumption. Each mixed term also equals the Lebesgue integral:
for example integrating $g(y)$ from $x$ to $1$ leaves a polynomial
in $x$ of degree at most $\deg f+\deg g+1\le M+1$, to which the
first marginal rule applies. Hence $C(f,g)=0$.
The same argument applies to the reverse chamber and to every pair
of polynomials within the degree budget. Since the measures have no
common atoms, their cumulative error functions have no simultaneous
jumps. It follows that
\[
 \int E_{1,f}\,dE_{2,g}=0,\qquad
 \int E_{2,g}\,dE_{1,f}=0
 \quad(\deg f+\deg g\le M).
\]

Let $h$ be an antiderivative of $p$. Stieltjes integration by parts
for $h E_{1,f}E_{2,g}$ has no common-jump term. Its boundary term is
zero: cumulative errors at $1$ vanish by marginal exactness, and at
$0$ any endpoint masses cannot occur simultaneously. Therefore
\[
 \int p E_{1,f}E_{2,g}
 =-\int E_{1,f}\,dE_{2,hg}
  -\int E_{2,g}\,dE_{1,hf}=0,
\]
because $\deg f+\deg g+\deg h\le M$. This proves~\eqref{orthogonal}.
For a finitely supported $\mu_i$, on the complement of its finitely
many atoms the derivative of $E_{i,f}$ is $-f$. Thus $E_{i,f}=0$
almost everywhere implies $f=0$, proving the dimension statement.
\end{proof}

\paragraph{Equivalent starting condition.}
Under marginal exactness, the chamber condition is equivalent to
$\int E_{1,f}\,dE_{2,g}=0$ for all
$\deg f+\deg g\le M$; expanding $\delta_1\otimes\delta_2$ proves
both implications. The Lebesgue orthogonality in~\eqref{orthogonal}
is a consequence, not asserted to be sufficient by itself.

\paragraph{Relevance and limitation.}
Two small real discrete variables together with $u$ independent uniforms
have scalar marginal exactness through degree $u$ and chamber exactness
through degree $u$ (the conditional-order kernels form a Bernstein basis
on each triangle). Thus the theorem applies with $M=u-1$, giving
\eqref{orthogonal} through total degree $u-2$. The extra marginal degree
in the theorem should not be inferred from chamber degree alone.
Consequently matching their scalar quadrature moments separately is
insufficient. The two cumulative-error spaces must also be orthogonal.
The reviewed polynomial construction achieves this by positive density
completion. No argument here reduces its counts to $O(M^2)$, nor rules
out that sharper possibility.
\end{document}
